\documentclass{article}
\usepackage[margin=1in, a4paper]{geometry}
\usepackage{amsmath, amssymb, amsfonts}
\usepackage{graphicx}
\usepackage{xcolor}
\usepackage{listings}
\usepackage{booktabs}
\usepackage[utf8]{inputenc}
\usepackage[T1]{fontenc}
\usepackage{hyperref}
\hypersetup{colorlinks=true, urlcolor=blue, linkcolor=blue}
\usepackage{enumitem}
\usepackage{float}

\definecolor{codegreen}{rgb}{0,0.6,0}
\definecolor{codegray}{rgb}{0.5,0.5,0.5}
\definecolor{codepurple}{rgb}{0.58,0,0.82}
\definecolor{backcolour}{rgb}{0.99,0.99,0.99}
% Professional color palette for academic publishing
\definecolor{myblue}{cmyk}{0.8,0.4,0,0.1}
\definecolor{mygreen}{cmyk}{0.7,0,0.5,0.2}
\definecolor{myorange}{cmyk}{0,0.5,0.8,0.1}
\definecolor{mygray}{cmyk}{0,0,0,0.4}
\definecolor{arrowcolor}{cmyk}{0.9,0.5,0,0.3}
\definecolor{nodecolor}{cmyk}{0.6,0.2,0,0.05}
\definecolor{framecolor}{cmyk}{0.4,0.1,0,0.1}

\lstdefinestyle{mystyle}{
    backgroundcolor=\color{backcolour},   
    commentstyle=\color{codegreen},
    keywordstyle=\color{magenta},
    numberstyle=\tiny\color{codegray},
    stringstyle=\color{codepurple},
    basicstyle=\ttfamily\footnotesize,
    breakatwhitespace=false,         
    breaklines=true,                 
    captionpos=b,                    
    keepspaces=true,                 
    numbers=left,                    
    numbersep=5pt,                  
    showspaces=false,                
    showstringspaces=false,
    showtabs=false,                  
    tabsize=2
}
\lstset{style=mystyle}

\title{The Logistics \& Supply Chain Problem-Solving Playbook\\
Chapter 82: The Single Facility Location Problem}
\author{Operations Research \& Supply Chain Analytics}
\date{\today}

\begin{document}

\maketitle

\section{The Single Facility Location Problem (Center of Gravity)}

The facility location problem represents one of the most fundamental challenges in supply chain design, where organizations must determine the optimal geographic position for a new distribution center, warehouse, or service facility. The center of gravity method provides an elegant mathematical approach to minimize transportation costs and distances by finding the weighted centroid of demand locations. This problem becomes particularly critical when companies expand their operations, consolidate facilities, or respond to changing market conditions where the optimal location can significantly impact operational costs, service levels, and competitive advantage.

\subsection{The Scenario: MegaCorp's Distribution Network Expansion}

MegaCorp, a major consumer electronics retailer, operates five retail stores across a metropolitan area and needs to establish a new distribution center to serve all locations efficiently. Each store has different weekly demand volumes, and the company wants to minimize the total transportation cost by optimizing the distribution center's location. The stores are located at coordinates (2,3), (8,1), (6,7), (12,4), and (4,9) with weekly demands of 150, 200, 180, 160, and 110 units respectively. Transportation cost is directly proportional to both distance and demand volume, making this a classic weighted facility location problem where the optimal position minimizes the sum of weighted distances to all demand points.

\subsection{Tier 1: The Pen \& Paper Method (Mathematical Formulation)}

The center of gravity method transforms the facility location problem into a weighted centroid calculation, where each demand location exerts a "gravitational pull" proportional to its volume. The mathematical foundation rests on minimizing the total weighted Euclidean distance from the new facility to all existing locations.

\textbf{Mathematical Model}

Let us define the following sets and parameters:
\begin{align}
I &= \{1, 2, \ldots, n\} \text{ (set of existing demand locations)} \\
(x_i, y_i) &= \text{coordinates of location } i \in I \\
w_i &= \text{weight (demand volume) at location } i \in I \\
(X, Y) &= \text{coordinates of the new facility}
\end{align}

The objective function minimizes the total weighted distance:
\min Z = \sum_{i \in I} w_i \sqrt{(X - x_i)^2 + (Y - y_i)^2}

For the center of gravity method, we use the assumption that distances are approximated linearly, leading to the weighted centroid formulation:
\begin{align}
X^* &= \frac{\sum_{i \in I} w_i \cdot x_i}{\sum_{i \in I} w_i} \\
Y^* &= \frac{\sum_{i \in I} w_i \cdot y_i}{\sum_{i \in I} w_i}
\end{align}

\textbf{Visual Representation}

\begin{figure}[htbp]
\centering
\includegraphics[width=\textwidth]{../figures-png/line82.fig1.png}
\caption{Center of Gravity Method Visualization}
\end{figure}

\textbf{Concrete Example Solution}

For MegaCorp's five stores, we calculate:

\textbf{Step 1: Calculate the X-coordinate}
\begin{align}
X^* &= \frac{150 \times 2 + 200 \times 8 + 180 \times 6 + 160 \times 12 + 110 \times 4}{150 + 200 + 180 + 160 + 110} \\
&= \frac{300 + 1600 + 1080 + 1920 + 440}{800} \\
&= \frac{5340}{800} = 6.675
\end{align}

\textbf{Step 2: Calculate the Y-coordinate}
\begin{align}
Y^* &= \frac{150 \times 3 + 200 \times 1 + 180 \times 7 + 160 \times 4 + 110 \times 9}{800} \\
&= \frac{450 + 200 + 1260 + 640 + 990}{800} \\
&= \frac{3540}{800} = 4.425
\end{align}

Therefore, the optimal location for MegaCorp's distribution center is at coordinates $(6.675, 4.425)$.

\textbf{Step 3: Verification of Total Weighted Distance}
Z^* = \sum_{i=1}^{5} w_i \sqrt{(6.675 - x_i)^2 + (4.425 - y_i)^2}

Computing each term:
\begin{align}
w_1 d_1 &= 150 \times \sqrt{(6.675-2)^2 + (4.425-3)^2} = 150 \times 4.89 = 734 \\
w_2 d_2 &= 200 \times \sqrt{(6.675-8)^2 + (4.425-1)^2} = 200 \times 3.67 = 734 \\
w_3 d_3 &= 180 \times \sqrt{(6.675-6)^2 + (4.425-7)^2} = 180 \times 2.63 = 473 \\
w_4 d_4 &= 160 \times \sqrt{(6.675-12)^2 + (4.425-4)^2} = 160 \times 5.34 = 854 \\
w_5 d_5 &= 110 \times \sqrt{(6.675-4)^2 + (4.425-9)^2} = 110 \times 5.20 = 572
\end{align}

Total weighted distance: $Z^* = 734 + 734 + 473 + 854 + 572 = 3367$ units.

\subsection{Tier 2: The Classic Heuristic (Geometric Iterative Refinement)}

While the center of gravity provides an excellent starting point, real-world facility location often requires consideration of discrete feasible locations, terrain constraints, and non-linear transportation costs. The geometric iterative refinement heuristic improves upon the basic centroid by incorporating distance corrections and feasibility constraints.

\textbf{Algorithm Overview}

The iterative refinement approach uses a multi-phase strategy: (1) compute the initial center of gravity, (2) evaluate a grid of candidate locations around this point, (3) apply feasibility filters, and (4) use local improvement to find the best discrete location.

\begin{figure}[htbp]
\centering
\includegraphics[width=\textwidth]{../figures-png/line82.fig2.png}
\caption{Iterative Grid Search Around Center of Gravity}
\end{figure}

\textbf{Python Implementation}

\lstinputlisting[language=Python,caption={Geometric Iterative Refinement Heuristic}]{.././codes-python/line82.py1.py}

\textbf{Concrete Example Output}

Running the heuristic on MegaCorp's data produces:

\begin{verbatim}
Center of Gravity: (6.675, 4.425)
Generated 49 candidate locations
Best grid candidate: (6.500, 4.500) with cost 3384.62
Final Results:
Optimal Location: (6.600, 4.400)
Total Weighted Distance: 3371.45
Average Distance per Unit: 4.214
\end{verbatim}

The heuristic finds a slightly different location than the pure mathematical center of gravity, accounting for discretization and feasibility constraints, with only a 0.13\% increase in total cost.

\subsection{Tier 3: The Advanced Algorithm (Particle Swarm Optimization)}

For complex facility location problems involving multiple constraints, non-linear cost functions, or multiple facilities, particle swarm optimization provides a robust metaheuristic approach. PSO models the solution space as a swarm of particles exploring potential facility locations, with each particle representing coordinates and learning from both personal and global best positions.

\textbf{PSO Adaptation for Facility Location}

Each particle $i$ in the swarm represents a potential facility location with:
\begin{align}
\mathbf{x}_i &= (x_i, y_i) \quad \text{(current position)} \\
\mathbf{v}_i &= (v_{x_i}, v_{y_i}) \quad \text{(velocity vector)} \\
\mathbf{p}_i &= (p_{x_i}, p_{y_i}) \quad \text{(personal best position)} \\
\mathbf{g} &= (g_x, g_y) \quad \text{(global best position)}
\end{align}

The velocity update equation incorporates both cognitive and social learning:
\[
\mathbf{v}_i^{t+1} = w \mathbf{v}_i^t + c_1 r_1 (\mathbf{p}_i - \mathbf{x}_i^t) + c_2 r_2 (\mathbf{g} - \mathbf{x}_i^t)
\]

Where $w$ is inertia weight, $c_1$ and $c_2$ are acceleration coefficients, and $r_1, r_2$ are random numbers in $[0,1]$.

\begin{figure}[htbp]
\centering
\includegraphics[width=\textwidth]{../figures-png/line82.fig3.png}
\caption{Particle Swarm Optimization Dynamics}
\end{figure}

\textbf{Enhanced PSO Implementation}

\lstinputlisting[language=Python,caption={Particle Swarm Optimization for Facility Location}]{.././codes-python/line82.py2.py}

\textbf{Concrete Example Output}

\begin{verbatim}
Starting PSO optimization...
Iteration 0: Best fitness = 4235.67
Iteration 20: Best fitness = 3387.45
Iteration 40: Best fitness = 3372.18
Iteration 60: Best fitness = 3368.91
Iteration 80: Best fitness = 3367.22
Iteration 100: Best fitness = 3366.85
Iteration 120: Best fitness = 3366.82
Iteration 140: Best fitness = 3366.81
Iteration 160: Best fitness = 3366.81
Iteration 180: Best fitness = 3366.81

PSO Results:
Optimal Location: (6.674, 4.427)
Total Weighted Distance: 3366.81

Convergence Analysis:
Total Improvement: 20.51%
Converged at Iteration: 85
Improvement over Center of Gravity: 0.01%
\end{verbatim}

The PSO approach demonstrates several advantages: it naturally handles constraints through penalty functions, explores the solution space more thoroughly than gradient-based methods, and provides insight into convergence behavior. The algorithm converged to a solution very close to the mathematical optimum while respecting practical constraints.

\section{Practice Questions}

\textbf{Tier 1 Questions (Mathematical Formulation)}

1. \textbf{Regional Distribution Center Planning:} A company has four manufacturing plants located at coordinates (5,2), (10,8), (15,3), and (7,12) with monthly production volumes of 500, 800, 650, and 400 tons respectively. Transportation cost is \$2 per ton per unit distance. Calculate the optimal distribution center location using the center of gravity method and determine the total monthly transportation cost.

2. \textbf{Multi-Product Facility Location:} A logistics company serves five cities with different products having varying transportation costs. City A at (3,4) requires 200 units at \$1.5/unit/distance, City B at (9,1) requires 300 units at \$2.0/unit/distance, City C at (6,8) requires 150 units at \$1.8/unit/distance, City D at (12,5) requires 250 units at \$1.6/unit/distance, and City E at (2,7) requires 180 units at \$2.2/unit/distance. Find the cost-weighted optimal facility location.

\textbf{Tier 2 Questions (Heuristic Implementation)}

3. \textbf{Constrained Grid Search:} Modify the geometric iterative refinement heuristic to handle a scenario where potential facility locations are restricted to a discrete set of 20 predetermined sites. The sites are located on a 4×5 grid with coordinates (i,j) where i ∈ \{2,4,6,8\} and j ∈ \{1,3,5,7,9\}. Three sites at (4,3), (6,7), and (8,5) are unavailable due to zoning restrictions. Implement the selection algorithm and test with the MegaCorp dataset.

4. \textbf{Multi-Objective Heuristic:} Design a heuristic that simultaneously minimizes transportation cost and maximizes service coverage (measured as the number of demand points within a 3-unit radius). Use the weighted-sum approach with weights 0.7 for cost and 0.3 for coverage. Test with six demand points: (1,1), (4,2), (7,8), (10,3), (5,6), (8,9) with demands 100, 150, 120, 180, 90, 160 respectively.

\textbf{Tier 3 Questions (Advanced Metaheuristics)}

5. \textbf{Multi-Swarm PSO Comparison:} Implement a multi-swarm PSO variant where three sub-swarms of 15 particles each optimize different aspects: cost minimization, constraint satisfaction, and solution robustness. Compare convergence rates and solution quality against the standard PSO for a 10-location facility problem with mixed integer constraints.

6. \textbf{Hybrid Genetic-PSO Algorithm:} Design a hybrid algorithm that uses genetic algorithm for global exploration in the first 50\% of iterations, then switches to PSO for local exploitation. The GA should use tournament selection, two-point crossover, and Gaussian mutation. Test on a complex scenario with 8 demand locations, 3 restricted zones, and non-linear distance cost function $f(d) = d^{1.3}$ to account for economies of scale in transportation.

\end{document}
